\documentclass[11pt]{article}
\usepackage[margin=1.1in]{geometry}
\usepackage{amsmath,amssymb,amsthm}
\usepackage{booktabs}
\usepackage[colorlinks=true,linkcolor=blue,citecolor=blue]{hyperref}

\newtheorem{theorem}{Theorem}[section]
\newtheorem{lemma}[theorem]{Lemma}
\newtheorem{proposition}[theorem]{Proposition}
\newtheorem{conjecture}[theorem]{Conjecture}
\newtheorem{corollary}[theorem]{Corollary}
\theoremstyle{definition}
\newtheorem{definition}[theorem]{Definition}
\newtheorem{remark}[theorem]{Remark}
\newtheorem{question}[theorem]{Question}

\newcommand{\ps}{\mathrm{ps}}
\newcommand{\cp}{\mathrm{cap}}

\title{Peg solitaire on unicyclic graphs:\\ crowns, cycle-critical graphs, and two atlases}
\author{Krishna Harish\\[2pt]
\normalsize Elkins High School, Missouri City, Texas}
\date{July 2026}

\begin{document}
\maketitle

\begin{abstract}
We study peg solitaire on graphs in the sense of Beeler and Hoilman.
Our starting point is a pair of exhaustive computational atlases: the
solvability status, for every starting hole, of all $205{,}001$ trees
on up to $18$ vertices and all $61{,}130$ unicyclic graphs on up to
$14$ vertices. Prior exhaustive data stopped at all graphs on $7$
vertices and freely solvable trees on $10$. The atlases lead to
several theorems and one new phenomenon. For the \emph{crown}
$C(k;p)$, a cycle with $p$ pendant vertices at a single vertex, we
determine the solvability threshold $p^*(k)$ exactly for all
$k\le 31$ and prove that the capacity of a cycle is bounded by an
absolute constant: $C(k;p)$ is unsolvable for every $k$ once
$p\ge 13$. The sufficiency half is proved for all $k$ via a
two-jump purge gadget and a new family of single-cluster caterpillar
lemmas, which also yield a position law: a cluster of $p$ pendants is
tolerated on a long path only at spine position $2p-1$ (even paths)
or $2(p-1)$ (odd paths). We then isolate the graphs responsible for
the failure of a natural reduction: call a solvable unicyclic graph
\emph{cycle-critical} if every spanning tree is unsolvable. There
are $1{,}839$ such graphs on at most $14$ vertices, and they include an
infinite two-parameter family, the double crowns $DC(k,d)$, for which
we prove a parity dichotomy:
$DC(k,d)$ with two pendant pairs at distance $d$ is freely reachable
from its spanning trees when $d$ is odd, and cycle-critical when $d$
is even (solvability proved for all $k\equiv 0 \pmod 4$; criticality
reduced to, and following from, a caterpillar lemma verified at all
positions through spine length $20$). We show that the number of pegs
a cycle edge can rescue is unbounded: the triangle with $q$
pendants at each vertex is solvable yet all its spanning trees strand
$\Theta(q)$ pegs, refuting a bound the small-graph data had
suggested. Finally we prove that all interaction in a
crown game is bounded by constants independent of $k$, measure the
Myhill--Nerode complexity of the reachable configuration spaces, and
explain why the eventual periodicity of $k\mapsto$ ``$C(k;p)$ is
solvable'', which all our data supports, resists the standard
regularity toolkit. We also correct a boundary case of a published
theorem on windmill graphs. All computations are exact and
reproducible from the accompanying repository.
\end{abstract}

\section{Introduction}

Peg solitaire on graphs was introduced by Beeler and Hoilman
\cite{BH11}, lifting the classical board game \cite{Beasley,BCG,Bell}
to arbitrary connected graphs. Pegs occupy all vertices of a
connected graph $G$ except one, the starting hole; a move takes a
path $u\,v\,w$ with pegs on $u,v$ and a hole on $w$, jumps the peg
from $u$ to $w$, and removes the peg on $v$. The graph is
\emph{solvable} if from some starting hole a sequence of jumps leaves
a single peg, and \emph{freely solvable} if this works from every
hole. The central open problem of the area, posed in \cite{BH11} and
repeated in the survey \cite{dWK23}, is to characterize the solvable
trees; progress has come one family at a time, through the census of
graphs on at most seven vertices \cite{BG12}, double stars
\cite{BH12}, trees of diameter four \cite{BW15}, caterpillars
\cite{BGH17}, and banana trees \cite{dWK20}, among others. A parallel
line studies graph operations and modified rules: Cartesian
\cite{KdWcart} and corona \cite{KdWcorona} products, line graphs
\cite{KdWline}, reversible \cite{ES} and merging \cite{EW} moves,
double-jump solitaire \cite{BGdj}, the fool's solitaire number and
other extremal questions \cite{BR12,BGext}, graphs of large maximum
degree \cite{Mat}, pagoda-function certificates \cite{Kpagoda}, and
an infinite-board analogue tied to Conway's soldiers \cite{Vito}.

This paper approaches the subject from the unicyclic side, where, to
our knowledge, no systematic study existed. Every claim below grew
out of an exhaustive, exactly verified dataset rather than the other
way around, and we have tried to keep the boundary between what is
proved, what is machine-verified (and therefore proved, when the
claim is finite), and what is conjectured as sharp as possible.

Our contributions:

\begin{enumerate}
\item \textbf{Two atlases} (Section~\ref{sec:atlas}). The solvability
classification of all trees on $4\le n\le 18$ vertices ($205{,}001$
trees; the prior frontier was $n\le 10$, and only for free
solvability \cite{BG13}) and of all unicyclic graphs on $n\le 14$
vertices ($61{,}130$ graphs). About $61\%$ of large trees are
solvable and about $24\%$ freely solvable; unicyclic graphs are far
more permissive, at roughly $85\%$ solvable. As a first dividend the
atlas caught a boundary error in a published theorem: the windmill
variant $W(1,1)$ of \cite{BH12} is claimed freely solvable and is
not (Proposition~\ref{prop:windmill}).
\item \textbf{Single-cluster caterpillars and the position law}
(Section~\ref{sec:caterpillar}). We compute the pendant capacity
$\cp(P_m,i)$, the largest cluster of pendants attachable at spine
position $i$ of an $m$-vertex path without destroying solvability,
for all $m\le 20$. Even paths have capacity zero at depth five and
beyond; odd paths obey an eventually periodic profile peaking at
position six. The solvable stable positions follow the law $i=2p-1$
(even $m$) and $i=2(p-1)$ (odd $m$). A two-jump \emph{purge gadget}
turns machine-found base cases into theorems for all lengths
(Theorem~\ref{thm:families}).
\item \textbf{Crowns} (Section~\ref{sec:crowns}). Write $C(k;p)$ for
a $k$-cycle with $p$ pendants at one vertex. We determine the
threshold $p^*(k)$, below which $C(k;p)$ is solvable and above which
it is not, exactly for all $k\le 31$
(Theorem~\ref{thm:crown}): $p^*=4$ for odd $13\le k\le 31$, $p^*=2$
for even $10\le k\le 30$, with ten exceptional small values. The
sufficiency half holds for all $k$ by inheritance from the
caterpillar families. For necessity we prove a golden-ratio weight
bound (Theorem~\ref{thm:uniform}): \emph{no} cycle, of any length,
tolerates $13$ pendants at a vertex. En route we prove that a
fully pegged path can deliver at most one peg to an adjacent drained
cell, ever (Lemma~\ref{lem:oneshot}), and that all interface
activity in a crown game is bounded by constants independent of $k$
(Section~\ref{sec:periodicity}).
\item \textbf{Cycle-critical graphs} (Section~\ref{sec:critical}). A
solvable unicyclic graph is cycle-critical if every spanning tree is
unsolvable: the cycle edge is load-bearing. The atlas contains
$1{,}839$ examples on $n\le 14$ vertices; the smallest is $C_4$ with
one pendant on each of two opposite vertices. Up to $n=14$ they fall
in two tiers (the best spanning tree strands two pegs in $1{,}810$
cases and three in the remaining $29$), but this does not persist:
the rescue $\min_T\ps(T)-\ps(G)$ is in fact unbounded
(Theorem~\ref{thm:rescue}), so cycle-critical graphs have arbitrarily
many tiers.
\item \textbf{Double crowns and a parity dichotomy}
(Section~\ref{sec:dc}). $DC(k,d)$ is a $k$-cycle with two pendants
at $x_0$ and two at $x_d$. For odd $d$ we prove $DC(k,d)$ solvable
together with its spanning trees (Theorem~\ref{thm:oddside}); for
even $d$ it is cycle-critical in every case we can test, and we
prove the solvability half for the antipodal family for all
$k\equiv 0\pmod 4$, $k\ge 16$, by an explicit circulating play
family (Theorem~\ref{thm:dcs}). The solvable starting holes obey a
clean law: they sit at distances $\{3,6,k/2-6,k/2-3\}$ from the
clusters (Proposition~\ref{prop:distlaw}).
\item \textbf{The periodicity question}
(Section~\ref{sec:periodicity}). All data says $k\mapsto p^*(k)$ is
eventually periodic with period two. We prove the three
bounded-interface lemmas that any closure argument will need, and
then measure why the obvious closure fails: the reachable
configuration languages of crown games have super-linearly growing
Myhill--Nerode complexity under the natural encodings, so no
$k$-uniform regular invariant exists. The question joins the
octal-games periodicity problem \cite{GuyNowakowski} in the class of
statements that are almost certainly true, cheap to verify
instance-by-instance, and apparently expensive to close.
\end{enumerate}

Throughout, ``machine-verified'' means an exhaustive, exact search
whose code and outputs are in the public repository accompanying this
paper. Finite claims verified this way are proved; we reserve
``conjecture'' for statements quantified over infinitely many
untested cases. Two of our own intermediate models contained errors
that were caught by cross-implementation disagreement and by testing
consequences against known instances; we record only the corrected
statements.

\section{Preliminaries}\label{sec:prelim}

Graphs are finite, simple, connected. $P_n$ and $C_n$ denote the
path and cycle on $n$ vertices. A \emph{pendant} is a vertex of
degree one. For a graph $G$ with starting hole $s$, a
\emph{terminal state} is one admitting no jump; $G$ is solvable from
$s$ if some play from $s$ terminates with one peg. We write
$\ps(G)$ for the minimum number of pegs over all terminal states and
all starting holes, so $G$ is solvable iff $\ps(G)=1$.

We use three known results freely.

\begin{proposition}[{\cite{BH11,BH12}}]\label{prop:known}
$P_n$ is solvable iff $n$ is even or $n=3$, and freely solvable iff
$n=2$. $C_n$ is solvable iff $n$ is even or $n=3$. The double star
$DS(L,R)$, $L\ge R$, is solvable iff $L\le R+1$.
\end{proposition}

\begin{proposition}[Inheritance {\cite{BGH17}}]\label{prop:inherit}
If a spanning subgraph of $G$ is solvable from hole $s$, so is $G$.
\end{proposition}

Our solver performs exact memoized search over peg configurations,
recording for every starting hole the minimum and maximum terminal
peg counts. It was validated against every exactly-stated published
theorem we could encode: $122$ instances covering paths, cycles,
complete graphs, stars, double stars and windmills, all of which
agree, with one exception that turned out to be an error in the
literature rather than in the solver.

\begin{proposition}\label{prop:windmill}
Theorem 2.2(ii) of \cite{BH12} states that the windmill variant
$W(P,B)$, a universal vertex $u$ carrying $B$ triangles and $P$
pendants, is freely solvable iff $P\le 2B-1$ and $(P,B)\ne(0,2)$.
This fails at $(P,B)=(1,1)$: $W(1,1)$ is solvable but not freely
solvable.
\end{proposition}

\begin{proof}
Start the hole at the pendant $p_1$. Every first jump must land in
$p_1$ and jump over $u$, since $u$ is the only neighbour of $p_1$.
Either blade vertex may jump; both leave the two remaining pegs on
$b_i$ and $p_1$, which are non-adjacent, so the game ends with two
pegs. The published sufficiency proof checks holes on $u$ and on
blades but not on pendants. (Solvability from blade holes is easily
checked, so the graph is solvable.)
\end{proof}

\section{The atlases}\label{sec:atlas}

Trees were generated by the Wright--Richmond--Odlyzko--McKay
algorithm; unicyclic graphs as necklaces of rooted trees under the
dihedral action, validated in two independent ways (counts match
OEIS A001429, and a Burnside/cycle-index computation matches per
cycle length). Every graph was solved exactly for every starting
hole.

\begin{proposition}\label{prop:atlas}
Of the $205{,}001$ trees with $4\le n\le 18$: at $n=18$, $61.0\%$
are solvable and $23.6\%$ freely solvable. There are no freely
solvable trees at $n\in\{4,5,7\}$. Of the $61{,}130$ unicyclic
graphs with $4\le n\le 14$: at $n=14$, $84.7\%$ are solvable and
more than half are freely solvable.
\end{proposition}

The contrast between $61\%$ and $85\%$ reflects the central theme of
the paper: a single extra edge changes the game substantially, and
the remaining sections examine how.

\section{Single-cluster caterpillars}\label{sec:caterpillar}

Write $P_m(i\!:\!p)$ for the caterpillar consisting of a spine
$x_0\cdots x_{m-1}$ with $p$ pendants at $x_i$, and define the
\emph{pendant capacity} $\cp(P_m,i)$ as the largest $p$ for which
$P_m(i\!:\!p)$ is solvable (from some hole), or $0$ if none.
Beeler, Green and Harper \cite{BGH17} studied caterpillars with
clusters at the spine ends; interior single clusters seem not to have
been treated.

\begin{proposition}[machine-verified, $m\le 20$]\label{prop:cap}
For even $m\ge 14$ the capacity profile in $i=0,1,2,\dots$ is
$0,1,1,2,2,0,0,\dots$: in particular a single pendant at depth
$\ge 5$ makes an even path unsolvable. For odd $m\ge 13$ the profile
is $2,0,2,0,3,0,4,0,2,0,2,\dots$: zero at odd depths, peaking at
$i=6$, settling to $2$. Small $m$ carry exceptions, e.g.\
$\cp(P_{11},4)=5$ and $\cp(P_9,2)=\cp(P_9,4)=4$.
\end{proposition}

The stable solvable positions obey a law: for a cluster of size $p$,
position $2p-1$ on even spines and $2(p-1)$ on odd spines. Each
pendant needs two spine vertices of runway. The following gadget
turns the finite data into theorems.

\begin{lemma}[Purge gadget]\label{lem:purge}
Let $H$ contain an induced path $y\,x_{k-2}\,x_{k-1}\,x_k\,x_{k+1}$
in which $x_{k-1},x_k,x_{k+1}$ have no further neighbours. From the
full configuration with hole at $x_k$, the two jumps
$x_{k-2}\!>\!x_{k-1}\!>\!x_k$ and $x_{k+1}\!>\!x_k\!>\!x_{k-1}$
leave $x_k,x_{k+1}$ permanently empty and the graph
$H-\{x_k,x_{k+1}\}$ fully pegged with hole at $x_{k-2}$. Hence if
$H-\{x_k,x_{k+1}\}$ is solvable from $x_{k-2}$, then $H$ is solvable
from $x_k$.
\end{lemma}

\begin{proof}
Legality of the two jumps is immediate. Afterwards no path of length
two reaches $x_k$ or $x_{k+1}$ through a pegged middle, and any play
of $H-\{x_k,x_{k+1}\}$ is a play of $H$ verbatim, since the two dead
vertices never regain pegs.
\end{proof}

\begin{theorem}\label{thm:families}
For all spine lengths of the appropriate parity,
$P_k(1\!:\!1)$ and $P_k(3\!:\!2)$ are solvable for even $k\ge 6$;
$P_k(0\!:\!1)$ and $P_k(2\!:\!2)$ for odd $k\ge 7$;
$P_k(4\!:\!3)$ for odd $k\ge 9$; and $P_k(6\!:\!4)$ for odd
$k\ge 13$. In each case the starting hole may be taken at
$x_{k-2}$.
\end{theorem}

\begin{proof}
Induction on $k$ in steps of two. The base cases are the following
explicit plays, each verified mechanically (spine vertices are
numbered $0,\dots,k-1$, pendants from $k$ upward):
\begin{center}\footnotesize
\begin{tabular}{ll}
\toprule
$P_6(1\!:\!1)$ & $2{>}3{>}4\ 0{>}1{>}2\ 5{>}4{>}3\ 3{>}2{>}1\ 6{>}1{>}0$\\
$P_6(3\!:\!2)$ & $6{>}3{>}4\ 1{>}2{>}3\ 7{>}3{>}2\ 5{>}4{>}3\ 3{>}2{>}1\ 0{>}1{>}2$\\
$P_7(0\!:\!1)$ & $3{>}4{>}5\ 1{>}2{>}3\ 7{>}0{>}1\ 6{>}5{>}4\ 4{>}3{>}2\ 2{>}1{>}0$\\
$P_7(2\!:\!2)$ & $3{>}4{>}5\ 1{>}2{>}3\ 6{>}5{>}4\ 4{>}3{>}2\ 7{>}2{>}1\ 0{>}1{>}2\ 8{>}2{>}1$\\
$P_9(4\!:\!3)$ & $5{>}6{>}7\ 3{>}4{>}5\ 8{>}7{>}6\ 6{>}5{>}4\ 9{>}4{>}3\ 2{>}3{>}4\ 0{>}1{>}2\ 10{>}4{>}3\ 2{>}3{>}4\ 11{>}4{>}3$\\
$P_{13}(6\!:\!4)$ & $9{>}10{>}11\ 7{>}8{>}9\ 13{>}6{>}7\ 4{>}5{>}6\ 2{>}3{>}4\ 0{>}1{>}2\ 14{>}6{>}5\ 5{>}4{>}3$\\
& $\quad 2{>}3{>}4\ 12{>}11{>}10\ 10{>}9{>}8\ 8{>}7{>}6\ 15{>}6{>}5\ 4{>}5{>}6\ 16{>}6{>}5$\\
\bottomrule
\end{tabular}
\end{center}
For the step, $P_{k+2}(i\!:\!p)$ satisfies the hypothesis of
Lemma~\ref{lem:purge} at its far end: the three outermost spine
vertices carry no pendants because $i\le 6 < k-2$. The gadget reduces
the full board with hole $x_k$ to $P_k(i\!:\!p)$ with hole
$x_{k-2}$, solvable by induction.
\end{proof}

\begin{remark}
Machine search shows these positions are the only stable ones: for
$p=2$ on long even spines the solvable positions are exactly
$i\in\{3,4\}$, and for $p=4$ on long odd spines exactly $i=6$.
Single-cluster caterpillars cap at $p=2$ (even spine) and $p=4$
(odd), which the next section shows are precisely the crown
capacities.
\end{remark}

\section{Crowns}\label{sec:crowns}

Let $C(k;p)$ be the cycle $x_0x_1\cdots x_{k-1}$ with $p\ge 1$
pendants at $x_0$.

\begin{theorem}\label{thm:crown}
For all $3\le k\le 31$, $C(k;p)$ is solvable iff $p\le p^*(k)$,
where $p^*(k)=4$ for odd $k\ge 13$, $p^*(k)=2$ for even $k\ge 10$,
and $(p^*(3),\dots,p^*(12))=(2,1,3,2,3,3,4,2,5,2)$. Moreover the
solvable range is an interval in $p$, and, for every $k$ in this
range, $p^*(k)=\max_i \cp(P_k,i)$: deleting the right cycle edge
never costs anything.
\end{theorem}

\begin{proof}
Sufficiency for large $k$ is Theorem~\ref{thm:families} plus
Proposition~\ref{prop:inherit}: delete the cycle edge at the stable
distance from $x_0$ and inherit from the caterpillar. Exceptional
rows and small cases are finite and machine-witnessed. Necessity for
$p^*<p\le 13$ was verified exhaustively, from every starting hole,
for even $k\le 30$ and odd $k\le 31$ ($211$ instances), using a
pendant-symmetric solver (pendants are interchangeable, so a
configuration is a cycle mask plus a counter) cross-validated against
the direct solver. Necessity for $p\ge 13$ and \emph{all} $k$ is
Theorem~\ref{thm:uniform} below.
\end{proof}

Before the uniform bound, two remarks. The theorem has a
counterintuitive corollary.

\begin{corollary}\label{cor:rescue-odd}
Bare odd cycles $C_k$, $k\ge 5$, are unsolvable, yet $C(k;p)$ is
solvable for $1\le p\le p^*(k)$: attaching a single pendant rescues
any odd cycle.
\end{corollary}

Also, for fixed $p$ the map $k\mapsto[\,p\le p^*(k)\,]$ appears
eventually periodic with period two; Section~\ref{sec:periodicity}
discusses why we do not know how to close this and state it here
only for $k\le 31$.

\subsection{A uniform capacity bound}

Let $\varphi$ be the golden ratio and $s=1/\varphi$, so
$s^2+s=1$.

\begin{theorem}\label{thm:uniform}
For every $k\ge 3$ and every $p\ge 13$, $C(k;p)$ is unsolvable.
\end{theorem}

\begin{proof}
Weight the cycle by $w(x_j)=s^{\min(j,k-j)}$ and give each pendant
weight $t\in(0,s^2)$. Let $W$ be the total weight of pegged
vertices. A jump along $u\,v\,w$ changes $W$ by $w(w)-w(u)-w(v)$.
Jumps toward $x_0$ along either arc change $W$ by
$s^d-s^{d+1}-s^{d+2}=s^d(1-s-s^2)=0$; jumps away, and the finitely
many antipodal orientations, are strictly negative (e.g.\ the
antipodal even case gives $-s^{k/2}$). Refills
$x_2\!>\!x_1\!>\!x_0$ change $W$ by $1-s-s^2=0$. Classify the jumps
that consume the peg on $x_0$: a pendant exiting over $x_0$ burns
$1+t-s>1-s=s^2$; a pendant hopping to an empty pendant slot burns
$1$; a cycle neighbour feeding a pendant burns $1+s-t>0$; $x_0$
jumping into the cycle burns $1+s-s^2>0$; a through jump
$x_1\!>\!x_0\!>\!x_{k-1}$ burns $1$. So no jump increases $W$, and
each pendant-exit event burns more than $s^2$.

Pendants have degree one, so a pendant peg leaves the cluster only by
exiting over $x_0$, and pendant-to-pendant hops do not change the
number of pegged pendants. If $a$ pendant-exits and $c$ feeds occur
and $\varepsilon\in\{0,1\}$ pendant pegs remain at the end, then
$a-c=p-\varepsilon$, so $a\ge p-1$. The initial weight satisfies
$W_0<1+pt+2\sum_{d\ge1}s^d=1+pt+2\varphi$, and $W$ remains positive,
so $s^2(p-1)\le s^2a<1+2\varphi+pt$. Letting $t\to0$ and using
$1/s^2=\varphi+1$,
\[
p-1\;\le\;(1+2\varphi)(\varphi+1)\;=\;5\varphi+3\;=\;11.09\ldots,
\]
whence $p\le 12$.
\end{proof}

The constant is far from the truth ($p^*\le5$ everywhere), but it is
uniform in $k$, which is what the structure of
Theorem~\ref{thm:crown} needs: the sharp threshold is then a
statement about the finite strip $p^*<p\le 12$ only.

\subsection{One-shot delivery}

The mechanism limiting $p^*$ is the following. Consider a
path $1,2,\dots,m$, fully pegged, adjacent at cell $1$ to an external
cell $0$ that an adversary may drain at any time; a \emph{delivery}
is a jump $2\!>\!1\!>\!0$.

\begin{lemma}[One-shot delivery]\label{lem:oneshot}
A fully pegged path delivers at most one peg, regardless of length:
after the first delivery no second delivery is possible. With drains
at both ends, at most one delivery per end.
\end{lemma}

\begin{proof}
After a delivery, cells $1,2$ are empty. We claim that once cells
$j,j+1$ are empty and all cells left of $j$ are empty, cell $j$ can
never again hold a peg; the lemma follows with $j=1$. A peg can
arrive at $j$ only by the jump $j{+}2\!>\!j{+}1\!>\!j$, which
requires a peg at $j+1$; a peg can arrive at $j+1$ only from
$j{+}3\!>\!j{+}2\!>\!j{+}1$, which empties $j+2$ and $j+3$ and
recreates the hypothesis two cells to the right. The induction is
well-founded because the path is finite.
\end{proof}

In a crown, all refills of $x_0$ beyond the two supplied by the arc
must therefore be financed by pendant exits landing on the gate
cells, so each refill is itself paid for by a pendant exit. The
exact bookkeeping is delicate (Section~\ref{sec:periodicity}), but
the following consequence of the same weighting technique is not.

\begin{lemma}\label{lem:refills}
In any solving play of any $C(k;p)$, the number $r$ of jumps landing
on $x_0$ satisfies $r\le 10$; consequently the total number of jumps
involving $x_0$ or a pendant is at most $21$.
\end{lemma}

\begin{proof}
Let $\Phi$ be the pegged weight of the arc $x_1,\dots,x_{k-1}$ under
$w(x_j)=s^{\min(j,k-j)}$. Refills remove $s+s^2=1$ from $\Phi$;
pendant exits add $s$; through jumps are neutral; all other jumps do
not increase $\Phi$. Hence $r\le\Phi_0+sa\le2\varphi+sa$. Every
jump involving $x_0$ or a pendant consumes the peg on $x_0$, whose
supply is at most $1+r$, so $a\le1+r$ and
$r(1-s)\le2\varphi+s$, giving $r\le(2\varphi+s)/s^2=10.09\ldots$
\end{proof}

\section{Cycle-critical graphs}\label{sec:critical}

The equality $p^*(k)=\max_i\cp(P_k,i)$ of Theorem~\ref{thm:crown}
suggests a general principle: perhaps a unicyclic graph is solvable
iff some spanning tree is. The atlas refutes this decisively.

\begin{definition}
A solvable unicyclic graph $G$ is \emph{cycle-critical} if $G-e$ is
unsolvable for every cycle edge $e$, i.e.\ every spanning tree of
$G$ is unsolvable.
\end{definition}

This is not the edge-critical notion of \cite{BGedge}, which concerns
\emph{adding} an edge to an unsolvable graph; here we \emph{delete} a
cycle edge from a solvable one.

\begin{proposition}[machine-verified]\label{prop:critical}
Exactly $1{,}839$ of the $52{,}021$ solvable unicyclic graphs on
$n\le 14$ vertices are cycle-critical, with counts
$1,0,6,9,12,91,126,398,1196$ for $n=6,\dots,14$. The smallest is
$C_4$ with one pendant on each of two opposite vertices. For
$1{,}810$ of them the best spanning tree strands exactly two pegs;
for the remaining $29$, every spanning tree strands at least three.
No solvable unicyclic graph on $n\le14$ vertices has all spanning
trees stranding four or more, though larger ones do
(Theorem~\ref{thm:rescue}).
\end{proposition}

Since a play that never uses a jump path through $e$ is verbatim a
play of $G-e$, criticality says every solution \emph{circulates}:
it uses every cycle edge. We tested two structural explanations and
both fail: the two stranded pegs of the best spanning tree rarely
sit at the endpoints of the deleted edge (5 cases of 28 at
$n\le10$), and there are critical graphs all of whose solutions
have zero net flow around the cycle (gross usage does not cancel;
net usage can). Criticality is a statement about gross traversal,
and its characterization is open.

We turn to the \emph{size} of the rescue,
$R(G)=\min_T\ps(T)-\ps(G)$ (the minimum over spanning trees). On
$n\le 12$ it is small ($R=0$ in $7437$ cases, $1$ in $421$, $2$ in
$13$, never more), which tempts one to conjecture a universal
bound. That temptation is wrong (Theorem~\ref{thm:rescue} below), but
one inequality is clean and we record it first.

One direction is provable and pins down where the difficulty lies.
Fix an optimal play $P$ of $G$ (ending with $r=\ps(G)$ pegs) and a
cycle edge $e$. Replay $P$ greedily on $T=G-e$: process its jumps in
order, keeping each that is legal in the current $T$-state and
skipping the rest.

\begin{lemma}[Skip identity]\label{lem:skip}
The greedy replay ends with exactly $r+s_e$ pegs, where $s_e$ is the
number of jumps it skips; and $s_e=0$ if $P$ never traverses $e$.
Hence $R(G)\le\min_e s_e$, and $R(G)\le 0$ whenever some optimal play
of $G$ omits a cycle edge.
\end{lemma}

\begin{proof}
Both $G$ under $P$ and $T$ under the replay lose exactly one peg per
executed jump; $G$ executes all $N$ jumps to reach $r$ pegs from
$n{-}1$, and $T$ executes $N-s_e$, reaching $(n{-}1)-(N-s_e)=r+s_e$.
Before the first skip the two states agree, so the first skipped jump
must be one that uses $e$; if none does, there are no skips.
\end{proof}

In particular a positive rescue forces every optimal play to traverse
every cycle edge, which for solvable $G$ is exactly cycle-criticality.
This bounds $R$ by skips but does not bound it absolutely, and
nothing does. Let $T(q)$ be the triangle with $q$ pendants
at each of its three vertices.

\begin{theorem}[Rescue is unbounded]\label{thm:rescue}
$\ps(T(q))=1$ for every $q$, while every spanning tree of $T(q)$ has
$\ps\ge q-3$. Hence $R(T(q))\ge q-4\to\infty$: no bound $R\le B$
holds. The smallest explicit counterexample to $R\le 2$ is $T(5)$
($n=18$, $\ps=1$, every tree $\ps=4$, so $R=3$).
\end{theorem}

\begin{proof}
Solvability: with the hole at a vertex $a$, the jumps
$b_q\!>\!b\!>\!a$, $c_q\!>\!c\!>\!b$, $a_q\!>\!a\!>\!c$ delete one
pendant from each vertex and restore the triangle to its initial
state, reducing $T(q)$ to $T(q-1)$; the induction ends at
$T(0)=K_3$, solved by $b\!>\!c\!>\!a$. Lower bound: a spanning tree
deletes a triangle edge, say $ab$, leaving the bipartite spine
$a\!-\!c\!-\!b$ with $q$ pendants at each vertex. Colour vertices by
the parity of their distance to $c$: the even class
$\{c\}\cup\{a\text{- and }b\text{-pendants}\}$ has $1+2q$ vertices,
the odd class $\{a,b\}\cup\{c\text{-pendants}\}$ has $2+q$. Every jump
lowers the peg count of its middle vertex's colour by exactly one;
the only non-leaf vertices, hence the only possible middles, are
$a,b,c$, of which just $c$ is even. So even-class pegs disappear only
through $c$-middle jumps, each of which empties $c$, and a peg returns
to $c$ only through an $a$- or $b$-middle jump. Counting gives
(even removed) $\le 1+$ (odd removed), which forces at least $q-3$
pegs in any terminal state. The triangle itself escapes this ledger
because an odd cycle is not bipartite.
\end{proof}

So the cycle edge is worth $\Theta(q)$ pegs, not a constant, and the
apparent bound on $n\le 12$ was an artefact of range. This gives a
second infinite family of cycle-critical graphs, $T(q)$ for $q\ge 3$
(every tree strands $q-1\ge 2$ pegs), triangle-based rather than
even-cycle-based, and it shows the strand-count of a cycle-critical
graph is itself unbounded: the two tiers seen up to $n=14$
(Proposition~\ref{prop:critical}) do not continue.

The mechanism behind criticality and the caterpillar necessity is the
same, and can be stated cleanly.

\begin{lemma}[Refill necessity]\label{lem:refill}
Let $v$ have exactly two pendant neighbours $p_1,p_2$. Any play that
removes the pegs on both $p_1$ and $p_2$ contains a jump landing on
$v$, occurring after a moment when $v$ is empty.
\end{lemma}

\begin{proof}
A leaf $p_r$ is removed only as the jumper of a jump over its unique
neighbour $v$ (it cannot be a middle, and it is never created). Such
a jump needs $v$ pegged immediately before and leaves $v$ empty
immediately after. Let $J_1$, then $J_2$, be the jumps removing
$p_1,p_2$. After $J_1$ the vertex $v$ is empty; just before $J_2$ it
is pegged; so some intervening jump lands on $v$, its middle a
non-pendant neighbour of $v$.
\end{proof}

In a double-cluster caterpillar this forces both $x_i$ and $x_j$ to
be refilled from the spine. The inter-cluster segment has
$j-i-1$ vertices, odd exactly when the separation $j-i$ is even; that
the two refills and the spine clearance cannot be reconciled on an
odd interior segment is what Proposition~\ref{prop:catnec} asserts and
what we cannot yet prove; it is the same obstruction, on a path, as
the crown refill ledger.

\section{Double crowns}\label{sec:dc}

Let $DC(k,d)$ be the $k$-cycle with two pendants at $x_0$ and two at
$x_d$, $2\le d\le k/2$. This family witnesses cycle-criticality
infinitely often, and the deciding parameter is the parity of $d$.

\begin{theorem}[Odd separation]\label{thm:oddside}
For $d\in\{3,5,7\}$ and every even $k\ge 2d+2$, the caterpillar
obtained from $DC(k,d)$ by deleting the cycle edge $x_{k-1}x_0$,
namely the double-cluster caterpillar with pairs at spine positions
$0$ and $d$, is solvable with hole at $x_{k-2}$. Consequently
$DC(k,d)$ is solvable and not cycle-critical.
\end{theorem}

\begin{proof}
The far spine end is three pendant-free vertices, so
Lemma~\ref{lem:purge} applies and reduces $k+2$ to $k$ exactly as in
Theorem~\ref{thm:families}; the base cases at $k=2d+2$ are
machine-witnessed. The same scheme works for any fixed odd $d$; we
state the cases whose base plays we have exhibited. Machine search
shows, additionally, that the cluster pair may sit at any even
offset $i$ with the same conclusion, and at no odd offset.
\end{proof}

\begin{theorem}[Even separation: solvability]\label{thm:dcs}
For every $k\equiv 0\pmod 4$ with $k\ge 16$, the antipodal double
crown $DC(k,k/2)$ is solvable with starting hole $x_3$.
\end{theorem}

\begin{proof}
Write $d=k/2$ and let the pendants at $x_0$ be $q_1,q_2$ and at
$x_d$ be $q_3,q_4$. We exhibit the play in six phases and prove each
legal; cells not mentioned retain their state.

\emph{Phase A}: $5{>}4{>}3$, $2{>}3{>}4$, $0{>}1{>}2$. Directly
legal from the start (hole $x_3$); afterwards the holes are exactly
$\{0,1,3,5\}$.

\emph{Phase B}: for $j=0,1,\dots,(d-8)/2$, the pair
$(2j{+}7)\,{>}\,(2j{+}6)\,{>}\,(2j{+}5)$ then
$(2j{+}4)\,{>}\,(2j{+}5)\,{>}\,(2j{+}6)$. We claim the invariant:
after pair $j$ the holes are $\{0,1\}\cup\{3,\dots,2j{+}5\}\cup
\{2j{+}7\}$, with pegs on $2$, on $2j{+}6$, and on everything from
$2j{+}8$ onward. For $j=0$ this is a direct check from the state
after A. Given the invariant at $j$, pair $j{+}1$ is legal: cells
$2j{+}9,2j{+}8$ are pegged and $2j{+}7$ is a hole for the first
jump, which fills $2j{+}7$ and empties $2j{+}8,2j{+}9$; then
$2j{+}6$ (pegged by the invariant) jumps over the just-filled
$2j{+}7$ into the just-emptied $2j{+}8$. The new state satisfies the
invariant at $j{+}1$. After the last pair the holes are
$\{0,1\}\cup\{3,\dots,d{-}3\}\cup\{d{-}1\}$ and the pegs on the low
side are exactly $x_2$ and $x_{d-2}$; cells $x_d,\dots,x_{k-1}$ are
untouched.

\emph{Phase C}: $(k{-}2){>}(k{-}1){>}0$, then
$m{>}(m{+}1){>}(m{+}2)$ for $m=k{-}4,k{-}6,\dots,d{+}4$. The first
jump is legal ($x_0$ is a hole from A) and empties
$x_{k-2},x_{k-1}$. Each subsequent step uses two untouched pegged
cells $m,m{+}1$ and the hole at $m{+}2$ left by the previous step;
it fills $m{+}2$ and empties $m,m{+}1$. After C the pegged cells
above $x_d$ are exactly $x_{d+1},x_{d+2},x_{d+3}$ (never touched)
and $x_{d+6},x_{d+8},\dots,x_{k-2}$, together with $x_0$.

\emph{Phase D}: $q{>}0{>}1$ with $q=q_1$, then $2{>}1{>}0$, then
$q{>}0{>}(k{-}1)$ with $q=q_2$. Legal in turn: $x_0$ pegged and
$x_1$ empty; then $x_2,x_1$ pegged, $x_0$ empty; then $x_0$ pegged
again and $x_{k-1}$ empty from C. The cluster at $x_0$ is now gone
and $x_{k-1}$ is pegged.

\emph{Phase E}: $m{>}(m{-}1){>}(m{-}2)$ for
$m=k{-}1,k{-}3,\dots,d{+}7$. Step $m$ uses the peg at $m$ (from D
when $m=k{-}1$, else from the previous step), the peg at $m{-}1$
(these are the cells $x_{d+6},\dots,x_{k-2}$ left by C), and the
hole at $m{-}2$ (emptied by C). After E the only pegs above $x_d$
are $x_{d+1},x_{d+2},x_{d+3}$ and $x_{d+5}$.

\emph{Phase F}: $q{>}d{>}(d{-}1)$, $(d{+}2){>}(d{+}1){>}d$,
$q{>}d{>}(d{+}1)$, $(d{-}2){>}(d{-}1){>}d$, $d{>}(d{+}1){>}(d{+}2)$,
$(d{+}2){>}(d{+}3){>}(d{+}4)$, $(d{+}5){>}(d{+}4){>}(d{+}3)$.
Entering F the pegs are exactly
$x_{d-2},x_d,x_{d+1},x_{d+2},x_{d+3},x_{d+5}$ plus the two pendants
at $x_d$; each jump is checked directly against the previous line,
and the final state is a single peg on $x_{d+3}$.

The phase lengths sum to $k+2$, which equals the number of vertices
minus two, as they must. For $k=16$ phases B, C, E consist of one
pair, two jumps, and one jump respectively, and the general
description specializes correctly; the play has also been verified
mechanically for $k=16,20,24,28,32,36,40$.
\end{proof}

The criticality half follows from a caterpillar necessity statement
and a clean reduction. Write $P_m(2\!@\!i,2\!@\!j)$ for the path on
$m$ vertices with two pendants at $x_i$ and two at $x_j$.

\begin{proposition}[Caterpillar necessity; machine-verified]
\label{prop:catnec}
Let $m$ be even, $14\le m\le 20$. For every $i<j$ with $j-i$ even,
$P_m(2\!@\!i,2\!@\!j)$ is unsolvable, at all positions. The parity
of $m$ matters: for odd $m$ even-separation solvable pairs occur
(e.g.\ several at $m=19$). We conjecture the statement for all even
$m$.
\end{proposition}

\begin{lemma}[Reduction]\label{lem:dcreduce}
For even $k$, every spanning tree of $DC(k,d)$ is a double-cluster
caterpillar of spine length $k$ whose two pairs are separated by
$d$ or $k-d$.
\end{lemma}

\begin{proof}
Deleting a cycle edge turns the $k$-cycle into a path on its $k$
vertices carrying the two pairs; the loaded vertices are joined the
short way (separation $d$) or the long way (separation $k-d$)
according to which arc holds the deleted edge.
\end{proof}

\begin{theorem}[Even separation: criticality]\label{thm:evencrit}
Let $k,d$ be even. Then $d$, $k-d$ and the spine length $k$ are all
even, so by Lemma~\ref{lem:dcreduce} every spanning tree of
$DC(k,d)$ is an even-spine, even-separation double-cluster
caterpillar. Granting Proposition~\ref{prop:catnec} at spine length
$k$, all spanning trees are unsolvable, so $DC(k,d)$, once solvable,
is cycle-critical. This is unconditional for even $k\le 20$, and
verified directly through $k=24$ for
\[
(k,d)\in\{(8,4),(12,4),(12,6),(16,8),(20,4),(20,10),(22,6),(22,10),
(24,8),(24,12)\}.
\]
For $d$ odd, $DC(k,d)$ is freely solvable with solvable spanning trees
in every tested case.
\end{theorem}

The dichotomy is therefore a theorem in a precise sense: the odd case
is proved (Theorem~\ref{thm:oddside}); the even case is proved modulo
extending the finite Proposition~\ref{prop:catnec} to all even $m$, a
$k$-uniform necessity of the same character as the crown strip
(Section~\ref{sec:periodicity}), together with solvability, which is
unconditional for the antipodal family (Theorem~\ref{thm:dcs}) and
machine-verified otherwise. The solvable holes carry unexpected
structure:

\begin{proposition}[machine-verified, $k\le 32$]\label{prop:distlaw}
For $k\equiv0\pmod4$, $16\le k\le32$, the solvable starting holes
of $DC(k,k/2)$ are exactly the cycle vertices at distance
$3$, $6$, $k/2-6$ or $k/2-3$ from one of the two loaded vertices.
($DC(12,6)$ is freely solvable; $DC(8,4)$ has its own small
pattern.)
\end{proposition}

The recurring $3$ and $6$ are the caterpillar runway constants of
Section~\ref{sec:caterpillar}, which we do not believe is a
coincidence, though we cannot yet derive one law from the other.

\section{Bounded interfaces and the periodicity question}
\label{sec:periodicity}

Everything we know says the following is true.

\begin{conjecture}\label{conj:periodic}
For each fixed $p$, the set of $k$ with $C(k;p)$ solvable is
eventually periodic in $k$ (with period $2$ and threshold $13$, on
the evidence).
\end{conjecture}

This would extend Theorem~\ref{thm:crown} from $k\le31$ to all $k$.
We record what we can prove toward it, and then quantify the
obstruction.

All interaction in a crown game is uniformly bounded. Jumps
landing on $x_0$ number at most $10$ and jumps meeting the pendant
cluster at most $21$ (Lemma~\ref{lem:refills}); and, weighting
distances from the cell antipodal to $x_0$ instead, the same
technique bounds the number of jumps whose path crosses the
antipodal cut by an absolute constant ($25$ with our loose
accounting). So for large $k$ a solving play decomposes into $O(1)$
interface events separated by episodes of pure one-dimensional
solitaire on segments. One-dimensional solitaire has a rational
reachability theory in the sense of Moore and Eppstein \cite{ME},
and one would like to conclude regularity, hence eventual
periodicity, by composition.

The obstruction is that our segments are not free-standing: they
begin fully pegged and are driven at their boundary by the bounded
event stream. We tested the natural closure route, regular model
checking: exhibit a regular invariant containing the reachable
configurations and excluding near-solved ones, and verify closure
under the move relation symbolically. That route is quantitatively
impossible here. Computing the exact reachable configuration sets of
$C(k;3)$ from all starting holes and measuring Myhill--Nerode
residual counts (lower bounds on any recognizing DFA):

\begin{center}\small
\begin{tabular}{lccccc}
\toprule
$k$ & 10 & 12 & 14 & 16 & 18\\
\midrule
linear encoding & 41 & 93 & 186 & 308 & ---\\
interleaved encoding & 43 & 97 & 197 & 346 & 479\\
\bottomrule
\end{tabular}
\end{center}

The counts grow super-linearly under both the around-the-cycle
encoding and the two-arms-from-the-gadget encoding: the reachable
languages are not $k$-uniformly regular, so no small regular
invariant exists in these encodings. This does not contradict
\cite{ME}, whose regularity concerns undriven paths, and it does not
make Conjecture~\ref{conj:periodic} doubtful; it locates the
difficulty. Either a different abstraction tames driven segments, or
the periodicity needs a direct combinatorial argument. The situation
is reminiscent of the octal-games periodicity conjecture
\cite{GuyNowakowski}: overwhelming instance evidence, bounded-looking
behaviour, no closure.

The caterpillar necessity of Proposition~\ref{prop:catnec} is a twin
of this question, and behaves identically. Its reachable
configuration sets, measured the same way for the even-separation
family as the spine grows, have Myhill--Nerode residual counts
$49, 94, 233, 296, 497$ at $m=8,10,12,14,16$: again super-linear,
again non-regular. We also checked the two most natural certificates
directly and both fail: the bipartite two-colouring gives a resource
count with no separation-parity constraint, and the exact modular
invariants $\sum_{\mathrm{pegs}} w \bmod q$ collapse to the trivial
weighting, because at a vertex carrying two pendants the three
cherries $p_1 x_i p_2$, $p_1 x_i x_{i\pm1}$ force $2w(x_i)\equiv 0$
and equal neighbour weights, which propagate along the Fibonacci
spine recurrence to kill every nonzero solution. So both of our $k$-uniform necessity statements, the crown strip and
the even-separation caterpillar, are hard for the same structural
reason, and neither yields to the standard invariant toolkit. Closing
either is, in our view, the main problem left open by this work.

\section{Open problems}

\begin{enumerate}
\item Prove Conjecture~\ref{conj:periodic}, even for $p=3$.
\item Characterize the cycle-critical unicyclic graphs. Theorem~\ref{thm:rescue}
shows the rescue $\min_T\ps(T)-\ps(G)$ is unbounded, so no constant
bound holds; what graph parameter controls it? For $T(q)$ the rescue
is $q-2$, linear in the per-vertex cluster size, and the follow-up
data (rescue notes in the repository) point to the balance of the
three clusters, $b+c-a$ for a loaded triangle, as the governing
quantity.
\item Extend Proposition~\ref{prop:catnec} to all even spine lengths:
an even-spine double-cluster caterpillar with even separation is
unsolvable. By Theorem~\ref{thm:evencrit} this closes the criticality
half of the double crown dichotomy. The obstruction is not the
bipartite two-colouring resource count, which we checked gives no
separation-parity constraint; a pagoda function in the sense of
\cite{Kpagoda} is the natural candidate.
\item Explain the Distance-$\{3,6\}$ law
(Proposition~\ref{prop:distlaw}) from the runway law of
Section~\ref{sec:caterpillar}.
\item Characterize solvable double-cluster caterpillars
$P_m(2\!@\!i,\,2\!@\!j)$; our data suggests parity of $j-i$
dominates, mirroring Theorem~\ref{thm:oddside}.
\item The 29 tier-two critical graphs (all spanning trees strand
three pegs) await any structural explanation; 23 of the 29 are
triangles carrying one heavy cluster and one chain.
\end{enumerate}

\section*{Reproducibility}

All solvers, generators, verification scripts, datasets, and the
witness plays quoted above are available in the repository
accompanying this paper. Every finite claim marked machine-verified
can be reproduced from a cold start in under an hour on a laptop.

\begin{thebibliography}{99}

\bibitem{BH11} R.~A. Beeler and D.~P. Hoilman, Peg solitaire on
graphs, \emph{Discrete Math.} 311 (2011), 2198--2202.

\bibitem{BH12} R.~A. Beeler and D.~P. Hoilman, Peg solitaire on the
windmill and the double star graphs, \emph{Australas. J. Combin.} 52
(2012), 127--134.

\bibitem{BG13} R.~A. Beeler and A.~D. Gray, Freely solvable graphs in
peg solitaire, \emph{ISRN Combinatorics} (2013), Article 605279.

\bibitem{BW15} R.~A. Beeler and C.~A. Walvoort, Peg solitaire on
trees with diameter four, \emph{Australas. J. Combin.} 63 (2015),
321--332.

\bibitem{BGH17} R.~A. Beeler, H.~Green and R.~T. Harper, Peg
solitaire on caterpillars, \emph{Integers} 17 (2017), \#G1.

\bibitem{dWK20} J.-H. de Wiljes and M.~Kreh, Peg solitaire on banana
trees, \emph{Bull. Inst. Combin. Appl.} 90 (2020), 63--86.

\bibitem{dWK23} J.-H. de Wiljes and M.~Kreh, Peg solitaire on graphs
--- a survey, \emph{Asian-Eur. J. Math.} 16 (2023), 2350059.

\bibitem{ME} C.~Moore and D.~Eppstein, One-dimensional peg solitaire,
and duotaire, in \emph{More Games of No Chance}, MSRI Publ. 42,
Cambridge Univ. Press, 2002, 341--350.

\bibitem{GuyNowakowski} R.~K. Guy and R.~J. Nowakowski, Unsolved
problems in combinatorial games, in \emph{Games of No Chance}, MSRI
Publ., Cambridge Univ. Press (various editions).

\bibitem{Beasley} J.~D. Beasley, \emph{The Ins and Outs of Peg
Solitaire}, Oxford University Press, 1985.

\bibitem{BCG} E.~R. Berlekamp, J.~H. Conway and R.~K. Guy,
\emph{Winning Ways for Your Mathematical Plays}, 2nd ed., A~K Peters,
2001--2004.

\bibitem{Bell} G.~I. Bell, Solving triangular peg solitaire,
\emph{J. Integer Seq.} 11 (2008), Article 08.4.8.

\bibitem{BG12} R.~A. Beeler and A.~D. Gray, Peg solitaire on graphs
with seven vertices or less, \emph{Congr. Numer.} 211 (2012),
151--159.

\bibitem{BR12} R.~A. Beeler and T.~K. Rodriguez, Fool's solitaire on
graphs, \emph{Involve} 5 (2012), no.~4, 473--480.

\bibitem{BGext} R.~A. Beeler and A.~D. Gray, Extremal results for peg
solitaire on graphs, \emph{Bull. Inst. Combin. Appl.} 77 (2016),
30--42.

\bibitem{BGdj} R.~A. Beeler and A.~D. Gray, Double jump peg solitaire
on graphs, \emph{Bull. Inst. Combin. Appl.} 91 (2021), 80--93.

\bibitem{KdWcart} M.~Kreh and J.-H. de Wiljes, Peg solitaire on
Cartesian products of graphs, \emph{Graphs Combin.} 37 (2021),
907--917.

\bibitem{KdWcorona} M.~Kreh and J.-H. de Wiljes, Peg solitaire on
corona products, \emph{Bull. Inst. Combin. Appl.} 96 (2022),
107--118.

\bibitem{KdWline} M.~Kreh and J.-H. de Wiljes, Peg solitaire on line
graphs, \emph{Trans. Comb.} 13 (2024), no.~3, 257--277.

\bibitem{Kpagoda} M.~Kreh, Pagoda functions for peg solitaire on
graphs, \emph{Discrete Appl. Math.} 358 (2024), 184--202.

\bibitem{BGedge} R.~A. Beeler and A.~D. Gray, Examples of edge
critical graphs in peg solitaire, in \emph{Combinatorics, Graph
Theory and Computing} (SEICCGTC 2020), Springer Proc. Math. Stat.
388, Springer, 2022, pp.~157--169.

\bibitem{ES} J.~Engbers and C.~Stocker, Reversible peg solitaire on
graphs, \emph{Discrete Math.} 338 (2015), 2014--2019.

\bibitem{EW} J.~Engbers and R.~Weber, Merging peg solitaire on graphs,
\emph{Involve} 11 (2018), 53--66.

\bibitem{Mat} N.~Matsumoto, Peg solitaire on graphs with large maximum
degree, \emph{Asian-Eur. J. Math.} 15 (2022), 2250057.

\bibitem{Vito} V.~Vito, Peg solitaire and Conway's soldiers on
infinite graphs, arXiv:2109.11128 (2021).

\end{thebibliography}

\end{document}
